\begin{exercise}[由衰减峰值反求串联电路参数]\label{ex:rev-c4-identify}
串联$RLC$回路与电源断开后自由振荡，已知$C$，电阻$R$、电感$L$未知。电容初始电压为$U_0>0$，电感初始电流为零。测得下一个同极性的电压峰值为$U_0/2$，出现于$t=T$。求$R,L$，并求这段时间内的总焦耳热。数值取$C=10\,\mu\mathrm F,T=2.00\,\mathrm{ms}$。
\ExpansionDiagram{fig:rev-c4-identify-q}{自由振荡回路：初始储能在电容中，回路不再接电源}{
\draw(0,0)--(0,2)--(1,2);\Res{1,2}{3,2}{R}\draw(3,2)--(4,2)--(4,1.15);\draw(3.7,1.15)--(4.3,1.15);\draw(3.7,.95)--(4.3,.95);\draw(4,.95)--(4,0)--(3,0);\node[right]at(4.35,1.05){$C$};
\draw(3,0)--(2.7,0);\draw(1.3,0)--(0,0);\foreach\x in{1.3,1.65,2,2.35}{\draw(\x,0)arc(180:0:.175);}\node[below]at(2,-.25){$L$};\node at(1.8,1){$U_C(0)=U_0$};\node at(1.8,.5){$i(0)=0$};\node[left]at(3.7,1.35){$+$};\node[left]at(3.7,.75){$-$};}
\end{exercise}
\begin{solution}
按图\ref{fig:rev-c4-identify-q}取电容原正极的电压为$U$。由回路方程与电容关系，
\[
 \ddot U+2\alpha\dot U+\omega_0^2U=0,\qquad
 \alpha=\frac{R}{2L},\quad\omega_0^2=\frac1{LC}.
\]
观察到同极性峰值重现，说明系统欠阻尼。令$\omega_d=\sqrt{\omega_0^2-\alpha^2}$，通解为$U=e^{-\alpha t}(A\cos\omega_dt+B\sin\omega_dt)$。由$U(0)=U_0$和$\dot U(0)=0$，先得$A=U_0$，再由$-\alpha A+\omega_dB=0$得$B=\alpha U_0/\omega_d$。因此
\[
 U=U_0e^{-\alpha t}\left(\cos\omega_dt+\frac\alpha{\omega_d}\sin\omega_dt\right),\quad
 \dot U=-\frac{\omega_0^2}{\omega_d}U_0e^{-\alpha t}\sin\omega_dt.
\]
\ExpansionDiagram{fig:rev-c4-identify-s}{同极性峰值间隔是完整阻尼周期，各电压极值处电流为零}{
\draw[axis](0,0)--(6,0)node[right]{$t/T$};\draw[axis](0,-1.9)--(0,2.5)node[above]{$U/U_0$};
\draw[curve,domain=0:1.1,samples=150]plot({5*\x},{2*exp(-ln(2)*\x)*(cos(360*\x)+ln(2)/(2*pi)*sin(360*\x))});
\draw[guide](5,0)--(5,1);\fill(0,2)circle(1.5pt);\fill(5,1)circle(1.5pt);\node[left]at(0,2){$1$};\node[above right]at(5,1){$1/2$};\node[below]at(5,0){$1$};\node[below]at(2.5,-1.45){反极性峰};}
由导数式，极值出现在$t=n\pi/\omega_d$；相邻同极性峰须取偶数$n$，所以图\ref{fig:rev-c4-identify-s}给
\[
 \omega_d=\frac{2\pi}{T},\qquad e^{-\alpha T}=\frac12,
 \qquad\alpha=\frac{\ln2}{T}.
\]
先恢复$\omega_0^2=\omega_d^2+\alpha^2$，再求元件参数：
\[
 \boxed{L=\frac{T^2}{C[4\pi^2+(\ln2)^2]},\qquad
 R=\frac{2T\ln2}{C[4\pi^2+(\ln2)^2]}}.
\]
代入数据得到$L\simeq10.01\,\mathrm{mH}$、$R\simeq6.94\,\Omega$。不能把测得的阻尼周期直接代入$T=2\pi\sqrt{LC}$；那会漏掉$\alpha^2$。

两个同极性峰处都有$\dot U=0$，故$i=0$，总储能分别为$CU_0^2/2$和$C(U_0/2)^2/2$。回路无外部供能，因此
\[
 Q_J=\frac12CU_0^2-\frac12C\left(\frac{U_0}2\right)^2
 =\frac38CU_0^2.
\]
电压峰值减半，对应能量变为原来的四分之一；若将相邻正、负峰之间的间隔误作$T$，还会把频率估计错一倍。
\end{solution}

\begin{exercise}[初始充电选择衰减模态]\label{ex:rev-c4-modes}
两个相同电容$C$分别接在节点$A,B$与公共地之间。每个节点各经一个电阻$R$接地，$A,B$之间另接一个电阻$R$。电源在$t=0$前已移去，初始节点电压为$V_A(0)=a,V_B(0)=b$。求两电压随时间的变化，说明怎样只激发慢模或快模；当$a=U,b=0$时，分别求三个电阻中产生的总热量。
\ExpansionDiagram{fig:rev-c4-modes-q}{电容与接地电阻并联，横向电阻只响应两节点电压差}{
\Res{0,2}{4,2}{R}\fill(0,2)circle(1.6pt)node[above left]{$A$};\fill(4,2)circle(1.6pt)node[above right]{$B$};
\foreach\x in{0,4}{\Res{\x,0}{\x,2}{R}\draw(\x,2)--({\x+1.3},2)--({\x+1.3},1.1);\draw({\x+1},1.1)--({\x+1.6},1.1);\draw({\x+1},.9)--({\x+1.6},.9);\draw({\x+1.3},.9)--({\x+1.3},0);\node[right]at({\x+1.6},1){$C$};}
\draw(0,0)--(5.3,0);\node[below]at(2.65,0){公共地，电位为零};}
\end{exercise}
\begin{solution}
图\ref{fig:rev-c4-modes-q}中，从节点$A$流出的电阻电流为$V_A/R+(V_A-V_B)/R$，使该节点电容电荷减少。两个节点方程为
\[
 C\dot V_A=-\frac{2V_A-V_B}{R},\qquad
 C\dot V_B=-\frac{2V_B-V_A}{R}.
\]
取$s=(V_A+V_B)/2$、$d=(V_A-V_B)/2$，方程相加、相减后分别得到
\[
 \dot s=-\frac{s}{RC},\qquad\dot d=-\frac{3d}{RC},\qquad
 s(0)=\frac{a+b}{2},\quad d(0)=\frac{a-b}{2}.
\]
两个模态分别乘积分因子，得
\[
 [s(t)e^{t/(RC)}]_0^T=0,\qquad[d(t)e^{3t/(RC)}]_0^T=0.
\]
因此$s(t)=s(0)e^{-t/(RC)}$、$d(t)=d(0)e^{-3t/(RC)}$；某个模态初始幅值为零时，它始终为零。用$V_A=s+d$、$V_B=s-d$还原，
\[
 \boxed{\begin{aligned}
 V_A&=\frac{a+b}{2}e^{-t/(RC)}+\frac{a-b}{2}e^{-3t/(RC)},\\
 V_B&=\frac{a+b}{2}e^{-t/(RC)}-\frac{a-b}{2}e^{-3t/(RC)}.
 \end{aligned}}
\]
\ExpansionDiagram{fig:rev-c4-modes-s}{同相模中横向无电流，反相模中横向电流加快放电}{
\begin{scope}\fill(0,1)circle(2pt)node[above]{$s$};\fill(2.4,1)circle(2pt)node[above]{$s$};\Res{0,1}{2.4,1}{R}\node at(1.2,.15){$I_{AB}=0$};\node at(1.2,-.55){时间常数$RC$};\end{scope}
\begin{scope}[xshift=5cm]\fill(0,1)circle(2pt)node[above]{$d$};\fill(2.4,1)circle(2pt)node[above]{$-d$};\Res{0,1}{2.4,1}{R}\draw[->](.4,.2)--(2,.2)node[midway,below]{$2d/R$};\node at(1.2,-.8){时间常数$RC/3$};\end{scope}}
图\ref{fig:rev-c4-modes-s}表明，$a=b$使$d(0)=0$，只留下慢模；$a=-b$使$s(0)=0$，只留下快模。观察到单一指数并不意味着电路只有一个独立储能变量，也可能是初态没有激发另一模态。

当$a=U,b=0$时，令$x=t/(RC)$，则$V_A=U(e^{-x}+e^{-3x})/2$、$V_B=U(e^{-x}-e^{-3x})/2$，且$\dd t/R=C\dd x$。
将电压平方展开后，只需积分三种指数。对任意$k>0$，
\[
 \int_0^\infty e^{-kx}\dd x=\lim_{X\to\infty}[-e^{-kx}/k]_0^X=\frac1k.
\]
两个接地电阻的平方展开分别为$U^2(e^{-2x}\pm2e^{-4x}+e^{-6x})/4$，故热量分别为
\begin{align*}
 Q_A&=\int_0^\infty\frac{V_A^2}{R}\dd t
 =\frac{CU^2}{4}\left(\frac12+\frac24+\frac16\right)
 =\frac7{24}CU^2,\\
 Q_B&=\int_0^\infty\frac{V_B^2}{R}\dd t
 =\frac{CU^2}{4}\left(\frac12-\frac24+\frac16\right)
 =\frac1{24}CU^2.
\end{align*}
横向电阻两端压差为$V_A-V_B=Ue^{-3x}$，故
\[
 Q_{AB}=CU^2\int_0^\infty e^{-6x}\dd x=\frac16CU^2.
\]
三者之和为$CU^2/2$，等于初始电容储能。慢、快模在总能量中正交；单个接地电阻的功率却包含两模态的交叉项，所以不能把总能量简单平均分给三个电阻。
\end{solution}
